\section*{第 \ref{ch:g7:lungs-and-gas-exchange} 章——肺与气体交换}

\begin{solution}{exo:g7:lungs-and-gas-exchange:1}
\omterm{def:g4:breathing:movements}{吸气}：肋间肌收缩，抬高并撑大胸廓，
\omterm{prop:g7:lungs-and-gas-exchange:ventilation}{膈肌}下降变平，被撑开的胸廓牵动
肺扩张，空气流入——这是\omterm{def:g3:bones-and-muscles:muscle}{肌肉}做功。
静息时\omterm{def:g4:breathing:movements}{呼气}：\omterm{def:g3:bones-and-muscles:muscle}{肌肉}舒张，被撑开的
胸廓弹回，把空气挤出——
几乎不耗能。
\end{solution}

\begin{solution}{exo:g7:lungs-and-gas-exchange:2}
肺下方的宽阔\omterm{def:g3:bones-and-muscles:muscle}{肌肉}层，封闭胸腔底部；
其收缩变平，
是每次\omterm{def:g4:breathing:movements}{吸气}中较深的那一部分。
\end{solution}

\begin{solution}{exo:g7:lungs-and-gas-exchange:3}
\omterm{def:g7:lungs-and-gas-exchange:alveoli}{肺泡}：数以亿计的薄壁
小囊泡，位于细支\omterm{def:g7:how-animals-breathe:four}{气管}末端，
每个都缠绕着细小\omterm{def:g4:heart-and-blood:vessels}{血管}。
\end{solution}

\begin{solution}{exo:g7:lungs-and-gas-exchange:4}
吸入气：\omterm{def:g7:what-is-respiration:respiration}{氧气}约占 $21$\%，几乎不含\omterm{def:g7:what-is-respiration:respiration}{二氧化碳}。
呼出气：\omterm{def:g7:what-is-respiration:respiration}{氧气}约占 $16$\%，
\omterm{def:g7:what-is-respiration:respiration}{二氧化碳}约占 $4$\%。两者之差——
吸入的\omterm{def:g7:what-is-respiration:respiration}{氧气}减少五个百分点、呼出的\omterm{def:g7:what-is-respiration:respiration}{二氧化碳}增加四个百分点——
正是气体交换本身。
\end{solution}

\begin{solution}{exo:g7:lungs-and-gas-exchange:5}
大面积、薄而湿润、内衬\omterm{def:g4:heart-and-blood:vessels}{血管}的交换面。与
鳃丝共用这一特征（
见\cref{prop:g7:how-animals-breathe:spec}）及
肠\omterm{prop:g7:digestion-and-nutrients:absorption}{绒毛}共用这一特征
（见\cref{ex:g7:digestion-and-nutrients:spec}）。
\end{solution}

\begin{solution}{exo:g7:lungs-and-gas-exchange:6}
焦油附着并刺激气道（超负荷的“扫帚”与咳嗽）；
\omterm{def:g7:lungs-and-gas-exchange:alveoli}{肺泡}壁破坏——交换面积丧失；
烟雾中的一种气体占据\omterm{prop:g4:journey-of-food:absorb}{血液}中的携氧位点；
气道变窄，扼制了“风箱”。
\end{solution}

\begin{solution}{exo:g7:lungs-and-gas-exchange:7}
\omterm{def:g4:breathing:movements}{吸气}时，须拉伸胸廓与肺，克服它们的弹性——
这一拉伸即是做功。静息时的\omterm{def:g4:breathing:movements}{呼气}，
便是弹性自行复原：扩张的胸廓与被拉伸的肺
回缩，把空气推出。
\end{solution}

\begin{solution}{exo:g7:lungs-and-gas-exchange:8}
气体交换既需要新鲜空气，也同样需要表面积：
交换只发生在空气与内衬\omterm{def:g4:heart-and-blood:vessels}{血管}的壁相接之处。
风箱加快运转，但\omterm{def:g7:lungs-and-gas-exchange:alveoli}{肺泡}数量更少、更膨大，
通过的却是已不存在的交换面——
缺失的因素正是交换面积本身。
\end{solution}

\begin{solution}{exo:g7:lungs-and-gas-exchange:9}
吸烟者的气道变窄（进气受阻），
\omterm{def:g7:lungs-and-gas-exchange:alveoli}{肺泡}表面积减少（每次\omterm{def:g7:what-is-respiration:respiration}{呼吸}的吸入量减少），
携氧位点被烟雾中的气体占据（每轮运送的\omterm{def:g7:what-is-respiration:respiration}{氧气}变少），
要承担同样的\omterm{def:g3:bones-and-muscles:muscle}{肌肉}耗氧，
\omterm{def:g4:heart-and-blood:heart}{心脏}就必须增加循环次数——因而
出现喘息，\omterm{prop:g4:heart-and-blood:pulse}{脉搏}达到冲刺时的水平。
\end{solution}

\begin{solution}{exo:g7:lungs-and-gas-exchange:10}
两者都是对\emph{呼出}时期的训练性控制：
在通常被动的弹性复原上，加入有节制的\omterm{def:g3:bones-and-muscles:muscle}{肌肉}用力——
潜水员要平稳地排出膨胀的空气，
吹奏者要把一次长长的呼出量控制得当。（他们的深\omterm{def:g4:breathing:movements}{吸气}
同样锻炼了\omterm{def:g4:breathing:movements}{吸气}时的用力。）
\end{solution}

\begin{solution}{exo:g7:lungs-and-gas-exchange:11}
更深、更快的\omterm{def:g7:what-is-respiration:respiration}{呼吸}，增大了\omterm{def:g7:lungs-and-gas-exchange:alveoli}{肺泡}端的
吸入量；更快、更有力的心跳，
增加了运输端的循环次数；改道使
\omterm{def:g3:bones-and-muscles:muscle}{肌肉}局部的\omterm{def:g4:heart-and-blood:vessels}{血管}在输送端扩张。三项调整，
共同倍增了一次输送的量——
（\cref{ex:g7:lungs-and-gas-exchange:accounting}中
收支平衡的账目）。
\end{solution}

\begin{solution}{exo:g7:lungs-and-gas-exchange:12}
边界：数以亿计的\omterm{def:g7:lungs-and-gas-exchange:alveoli}{肺泡}合计约
半个网球场大小的交换面积，折叠于
胸腔之中——每次\omterm{def:g7:what-is-respiration:respiration}{呼吸}里，五个百分点的\omterm{def:g7:what-is-respiration:respiration}{氧气}
向内穿过，四个百分点的\omterm{def:g7:what-is-respiration:respiration}{二氧化碳}
向外穿过。缓慢的火：烟雾中的焦油与刺激物
逐年侵蚀这一边界——包覆其
通路、塌陷其\omterm{def:g7:lungs-and-gas-exchange:alveoli}{肺泡}、夺取其运输
位点——这一损害是累积的、悄无声息的，且只能部分逆转，
好比边界的损毁快于重建。
\end{solution}

\begin{solution}{pb:g7:lungs-and-gas-exchange:1}
\textbf{1.} \omterm{def:g4:breathing:movements}{吸气}时期：肋间肌收缩，抬高并撑大胸廓，
\omterm{prop:g7:lungs-and-gas-exchange:ventilation}{膈肌}变平——需耗能。\omterm{def:g4:breathing:movements}{呼气}时期（静息时）：
\omterm{def:g3:bones-and-muscles:muscle}{肌肉}舒张，胸廓与肺弹回——
几乎不耗能。每分钟约循环 $15$ 次。

\textbf{2.} 肺附着于胸腔内壁：
扩张胸廓，肺也随之被拉开，
外界空气流入，充满被拉伸的空间。
这一“泵”的工作方式是改变空间的形状，
而非挤压器官。

\textbf{3.} 静息时每分钟约循环 $15$ 次；用力时
循环加快\emph{且}加深——\omterm{def:g7:what-is-respiration:respiration}{呼吸}次数更多，
每次吸入的空气也更多。

\textbf{4.} 喘息：快而浅的循环，几乎全靠
\omterm{prop:g7:lungs-and-gas-exchange:ventilation}{膈肌}完成。叹息：一次额外深长的扩张，随后
缓慢呼出——相当于弹簧复位。吹蜡烛：\omterm{def:g4:breathing:movements}{呼气}
阶段，\omterm{def:g3:bones-and-muscles:muscle}{肌肉}为弹簧添力。

\textbf{5.} \omterm{def:g7:lungs-and-gas-exchange:alveoli}{肺泡}：数以亿计的
薄壁小泡，每个外面都包裹着细\omterm{def:g4:heart-and-blood:vessels}{血管}，
合计面积约为半个网球场。

\textbf{6.} 面积大：即那半个网球场。薄：泡壁
仅一个\omterm{def:g5:first-look-at-cells:cell}{细胞}的宽度。湿润：其内衬。有\omterm{def:g4:heart-and-blood:vessels}{血管}包绕：
即外包的\omterm{def:g4:heart-and-blood:vessels}{血管}。四项要求全部满足。

\textbf{7.} $21$ 个百分点中留下了 5 个百分点——约
每次\omterm{def:g7:what-is-respiration:respiration}{呼吸}吸入的\omterm{def:g7:what-is-respiration:respiration}{氧气}约有四分之一被摄取。

\textbf{8.} 关于表面的湿润：呼出的空气
掠过潮湿（也较温暖）的壁面，因此呼出气
带走水分与热量——镜面起雾，
即表面维持的成本，每次\omterm{def:g7:what-is-respiration:respiration}{呼吸}都要支付。

\textbf{9.} 吸入：焦油附着并刺激气道。
泵：气道变窄，每次循环都受节流。表面：
\omterm{def:g7:lungs-and-gas-exchange:alveoli}{肺泡}壁破坏并相互融合——面积丧失。
运输：烟雾中的气体占据了\omterm{prop:g4:journey-of-food:absorb}{血液}的
携氧位点。

\textbf{10.} 咳嗽与感冒使气道发炎并清除——
机器修复后重新投入使用。\omterm{def:g7:lungs-and-gas-exchange:alveoli}{肺泡}壁一旦破损
不能重建：每个融合的\omterm{def:g7:lungs-and-gas-exchange:alveoli}{肺泡}囊都是永久丧失的
交换面积，属于资本性损失，而非
日常运营成本。

\textbf{11.} (1) 不让烟雾进入：资本性损失
永久存在。(2) \omterm{def:g7:what-is-respiration:respiration}{呼吸}洁净空气：过滤器与纤毛
的容量有限。(3) 负荷下锻炼“泵”：深
\omterm{def:g7:what-is-respiration:respiration}{呼吸}用力可维持\omterm{def:g3:bones-and-muscles:muscle}{肌肉}，并使整个
表面通气。(4) 寒冷多尘时经鼻吸入：
预处理可保护这套机器。

\textbf{12.} 在人体的诸多“泵”中，唯独这具“风箱”
既受自动调节支配，也受意志支配——
因为排出\omterm{def:g7:what-is-respiration:respiration}{二氧化碳}的同一股气流，
也是言语与歌唱的气流，而借给语言的机器
必须接受说话者的手来操控。
\end{solution}
